% Options for packages loaded elsewhere
\PassOptionsToPackage{unicode}{hyperref}
\PassOptionsToPackage{hyphens}{url}
\documentclass[
]{article}
\usepackage{xcolor}
\usepackage[margin=1in]{geometry}
\usepackage{amsmath,amssymb}
\setcounter{secnumdepth}{-\maxdimen} % remove section numbering
\usepackage{iftex}
\ifPDFTeX
  \usepackage[T1]{fontenc}
  \usepackage[utf8]{inputenc}
  \usepackage{textcomp} % provide euro and other symbols
\else % if luatex or xetex
  \usepackage{unicode-math} % this also loads fontspec
  \defaultfontfeatures{Scale=MatchLowercase}
  \defaultfontfeatures[\rmfamily]{Ligatures=TeX,Scale=1}
\fi
\usepackage{lmodern}
\ifPDFTeX\else
  % xetex/luatex font selection
\fi
% Use upquote if available, for straight quotes in verbatim environments
\IfFileExists{upquote.sty}{\usepackage{upquote}}{}
\IfFileExists{microtype.sty}{% use microtype if available
  \usepackage[]{microtype}
  \UseMicrotypeSet[protrusion]{basicmath} % disable protrusion for tt fonts
}{}
\makeatletter
\@ifundefined{KOMAClassName}{% if non-KOMA class
  \IfFileExists{parskip.sty}{%
    \usepackage{parskip}
  }{% else
    \setlength{\parindent}{0pt}
    \setlength{\parskip}{6pt plus 2pt minus 1pt}}
}{% if KOMA class
  \KOMAoptions{parskip=half}}
\makeatother
\usepackage{graphicx}
\makeatletter
\newsavebox\pandoc@box
\newcommand*\pandocbounded[1]{% scales image to fit in text height/width
  \sbox\pandoc@box{#1}%
  \Gscale@div\@tempa{\textheight}{\dimexpr\ht\pandoc@box+\dp\pandoc@box\relax}%
  \Gscale@div\@tempb{\linewidth}{\wd\pandoc@box}%
  \ifdim\@tempb\p@<\@tempa\p@\let\@tempa\@tempb\fi% select the smaller of both
  \ifdim\@tempa\p@<\p@\scalebox{\@tempa}{\usebox\pandoc@box}%
  \else\usebox{\pandoc@box}%
  \fi%
}
% Set default figure placement to htbp
\def\fps@figure{htbp}
\makeatother
\setlength{\emergencystretch}{3em} % prevent overfull lines
\providecommand{\tightlist}{%
  \setlength{\itemsep}{0pt}\setlength{\parskip}{0pt}}
\usepackage{multirow}
\usepackage{multicol}
\usepackage{colortbl}
\usepackage{hhline}
\newlength\Oldarrayrulewidth
\newlength\Oldtabcolsep
\usepackage{longtable}
\usepackage{array}
\usepackage{hyperref}
\usepackage{float}
\usepackage{wrapfig}
\usepackage{bookmark}
\IfFileExists{xurl.sty}{\usepackage{xurl}}{} % add URL line breaks if available
\urlstyle{same}
\hypersetup{
  pdftitle={Representatividade Indígena na Administração Pública Federal},
  pdfauthor={CGINF/DIGID/SGP/MGI},
  hidelinks,
  pdfcreator={LaTeX via pandoc}}

\title{Representatividade Indígena na Administração Pública Federal}
\author{CGINF/DIGID/SGP/MGI}
\date{24 abril, 2026}

\begin{document}
\maketitle

{
\setcounter{tocdepth}{2}
\tableofcontents
}
\section{Análise da
Representatividade}\label{anuxe1lise-da-representatividade}

Em Fevereiro/2026, a Administração Pública Federal (APF) contava com um
total de 2.247 vínculos ativos de pessoas que se autodeclaram como
indígenas, o que representa 0,40\% do total de ativos, conforme
representado na Tabela a seguir.

\global\setlength{\Oldarrayrulewidth}{\arrayrulewidth}

\global\setlength{\Oldtabcolsep}{\tabcolsep}

\setlength{\tabcolsep}{2pt}

\renewcommand*{\arraystretch}{1.5}



\providecommand{\ascline}[3]{\noalign{\global\arrayrulewidth #1}\arrayrulecolor[HTML]{#2}\cline{#3}}

\begin{longtable}[c]{|p{1.50in}|p{1.50in}|p{1.50in}}



\hhline{>{\arrayrulecolor[HTML]{666666}\global\arrayrulewidth=1.5pt}->{\arrayrulecolor[HTML]{666666}\global\arrayrulewidth=1.5pt}->{\arrayrulecolor[HTML]{666666}\global\arrayrulewidth=1.5pt}-}

\multicolumn{1}{>{\cellcolor[HTML]{0000AA}\raggedright}m{\dimexpr 1.5in+0\tabcolsep}}{\textcolor[HTML]{FFFFFF}{\fontsize{11}{11}\selectfont{Cor/Origem\ Étnica}}} & \multicolumn{1}{>{\cellcolor[HTML]{0000AA}\raggedleft}m{\dimexpr 1.5in+0\tabcolsep}}{\textcolor[HTML]{FFFFFF}{\fontsize{11}{11}\selectfont{Total\ de\ Vínculos}}} & \multicolumn{1}{>{\cellcolor[HTML]{0000AA}\raggedleft}m{\dimexpr 1.5in+0\tabcolsep}}{\textcolor[HTML]{FFFFFF}{\fontsize{11}{11}\selectfont{Participação\ (\%)}}} \\

\noalign{\global\arrayrulewidth 0pt}\arrayrulecolor[HTML]{000000}

\hhline{>{\arrayrulecolor[HTML]{666666}\global\arrayrulewidth=1.5pt}->{\arrayrulecolor[HTML]{666666}\global\arrayrulewidth=1.5pt}->{\arrayrulecolor[HTML]{666666}\global\arrayrulewidth=1.5pt}-}\endfirsthead 

\hhline{>{\arrayrulecolor[HTML]{666666}\global\arrayrulewidth=1.5pt}->{\arrayrulecolor[HTML]{666666}\global\arrayrulewidth=1.5pt}->{\arrayrulecolor[HTML]{666666}\global\arrayrulewidth=1.5pt}-}

\multicolumn{1}{>{\cellcolor[HTML]{0000AA}\raggedright}m{\dimexpr 1.5in+0\tabcolsep}}{\textcolor[HTML]{FFFFFF}{\fontsize{11}{11}\selectfont{Cor/Origem\ Étnica}}} & \multicolumn{1}{>{\cellcolor[HTML]{0000AA}\raggedleft}m{\dimexpr 1.5in+0\tabcolsep}}{\textcolor[HTML]{FFFFFF}{\fontsize{11}{11}\selectfont{Total\ de\ Vínculos}}} & \multicolumn{1}{>{\cellcolor[HTML]{0000AA}\raggedleft}m{\dimexpr 1.5in+0\tabcolsep}}{\textcolor[HTML]{FFFFFF}{\fontsize{11}{11}\selectfont{Participação\ (\%)}}} \\

\noalign{\global\arrayrulewidth 0pt}\arrayrulecolor[HTML]{000000}

\hhline{>{\arrayrulecolor[HTML]{666666}\global\arrayrulewidth=1.5pt}->{\arrayrulecolor[HTML]{666666}\global\arrayrulewidth=1.5pt}->{\arrayrulecolor[HTML]{666666}\global\arrayrulewidth=1.5pt}-}\endhead



\multicolumn{1}{>{\raggedright}m{\dimexpr 1.5in+0\tabcolsep}}{\textcolor[HTML]{000000}{\fontsize{11}{11}\selectfont{BRANCA}}} & \multicolumn{1}{>{\raggedleft}m{\dimexpr 1.5in+0\tabcolsep}}{\textcolor[HTML]{000000}{\fontsize{11}{11}\selectfont{320.441}}} & \multicolumn{1}{>{\raggedleft}m{\dimexpr 1.5in+0\tabcolsep}}{\textcolor[HTML]{000000}{\fontsize{11}{11}\selectfont{56,6\%}}} \\

\noalign{\global\arrayrulewidth 0pt}\arrayrulecolor[HTML]{000000}





\multicolumn{1}{>{\raggedright}m{\dimexpr 1.5in+0\tabcolsep}}{\textcolor[HTML]{000000}{\fontsize{11}{11}\selectfont{PARDA}}} & \multicolumn{1}{>{\raggedleft}m{\dimexpr 1.5in+0\tabcolsep}}{\textcolor[HTML]{000000}{\fontsize{11}{11}\selectfont{190.749}}} & \multicolumn{1}{>{\raggedleft}m{\dimexpr 1.5in+0\tabcolsep}}{\textcolor[HTML]{000000}{\fontsize{11}{11}\selectfont{33,7\%}}} \\

\noalign{\global\arrayrulewidth 0pt}\arrayrulecolor[HTML]{000000}





\multicolumn{1}{>{\raggedright}m{\dimexpr 1.5in+0\tabcolsep}}{\textcolor[HTML]{000000}{\fontsize{11}{11}\selectfont{PRETA}}} & \multicolumn{1}{>{\raggedleft}m{\dimexpr 1.5in+0\tabcolsep}}{\textcolor[HTML]{000000}{\fontsize{11}{11}\selectfont{40.413}}} & \multicolumn{1}{>{\raggedleft}m{\dimexpr 1.5in+0\tabcolsep}}{\textcolor[HTML]{000000}{\fontsize{11}{11}\selectfont{7,1\%}}} \\

\noalign{\global\arrayrulewidth 0pt}\arrayrulecolor[HTML]{000000}





\multicolumn{1}{>{\raggedright}m{\dimexpr 1.5in+0\tabcolsep}}{\textcolor[HTML]{000000}{\fontsize{11}{11}\selectfont{AMARELA}}} & \multicolumn{1}{>{\raggedleft}m{\dimexpr 1.5in+0\tabcolsep}}{\textcolor[HTML]{000000}{\fontsize{11}{11}\selectfont{11.787}}} & \multicolumn{1}{>{\raggedleft}m{\dimexpr 1.5in+0\tabcolsep}}{\textcolor[HTML]{000000}{\fontsize{11}{11}\selectfont{2,1\%}}} \\

\noalign{\global\arrayrulewidth 0pt}\arrayrulecolor[HTML]{000000}





\multicolumn{1}{>{\cellcolor[HTML]{FF7800}\raggedright}m{\dimexpr 1.5in+0\tabcolsep}}{\textcolor[HTML]{000000}{\fontsize{11}{11}\selectfont{INDIGENA}}} & \multicolumn{1}{>{\cellcolor[HTML]{FF7800}\raggedleft}m{\dimexpr 1.5in+0\tabcolsep}}{\textcolor[HTML]{000000}{\fontsize{11}{11}\selectfont{2.247}}} & \multicolumn{1}{>{\cellcolor[HTML]{FF7800}\raggedleft}m{\dimexpr 1.5in+0\tabcolsep}}{\textcolor[HTML]{000000}{\fontsize{11}{11}\selectfont{0,4\%}}} \\

\noalign{\global\arrayrulewidth 0pt}\arrayrulecolor[HTML]{000000}





\multicolumn{1}{>{\raggedright}m{\dimexpr 1.5in+0\tabcolsep}}{\textcolor[HTML]{000000}{\fontsize{11}{11}\selectfont{NAO\ INFORMADO}}} & \multicolumn{1}{>{\raggedleft}m{\dimexpr 1.5in+0\tabcolsep}}{\textcolor[HTML]{000000}{\fontsize{11}{11}\selectfont{50}}} & \multicolumn{1}{>{\raggedleft}m{\dimexpr 1.5in+0\tabcolsep}}{\textcolor[HTML]{000000}{\fontsize{11}{11}\selectfont{0,0\%}}} \\

\noalign{\global\arrayrulewidth 0pt}\arrayrulecolor[HTML]{000000}

\hhline{>{\arrayrulecolor[HTML]{666666}\global\arrayrulewidth=1.5pt}->{\arrayrulecolor[HTML]{666666}\global\arrayrulewidth=1.5pt}->{\arrayrulecolor[HTML]{666666}\global\arrayrulewidth=1.5pt}-}



\end{longtable}



\arrayrulecolor[HTML]{000000}

\global\setlength{\arrayrulewidth}{\Oldarrayrulewidth}

\global\setlength{\tabcolsep}{\Oldtabcolsep}

\renewcommand*{\arraystretch}{1}

Desde 2016, a representatividade de indígnas na APF alcançou seu ponto
mais baixo em 2019, com 1.983 vínculos (0,33\% do total), sendo que o
pico foi observado em 2023, com 2.411 vínculos (0,42\%).

Entre 2021 e 2022, é possível notar um aumento no total absoluto
acompanhado de uma queda no percentual, o que pode ser explicado por um
aumento no número total de vínculos mais acelerado para as outras
categorias de cor ou origem étnica.

\pandocbounded{\includegraphics[keepaspectratio]{index_files/figure-latex/grafico_serie_historica-1.png}}

As Unidades Federativas (UFs) com maior presença de indígenas nos órgãos
da APF são RR, AM e MT. Em cada uma destas UF, 3.39\%, 2.28\% e 1.17\%
dos vínculos, respectivamente, são exercidos por pessas desta etnia.

\pandocbounded{\includegraphics[keepaspectratio]{index_files/figure-latex/plot_uf-1.png}}

Ao analisar a representatividade de indígenas em cada órgão, a tabela a
seguir aponta Ministério Dos Povos Indígenas (MPI), com 33.03\% e
Fundação Nacional Dos Povos Indígenas (FUNAI), com 23.35\% como aqueles
em que há maior presença de indígenas com relação ao total de vínculos.
Juntos, estes órgãos contam com 2.156 vínculos, sendo que 23,8\% são de
pessoas indígenas. Por outro lado, estes mesmo órgãos concentram 22,9\%
dos vínculos ativos de indígenas.

\global\setlength{\Oldarrayrulewidth}{\arrayrulewidth}

\global\setlength{\Oldtabcolsep}{\tabcolsep}

\setlength{\tabcolsep}{2pt}

\renewcommand*{\arraystretch}{1.5}



\providecommand{\ascline}[3]{\noalign{\global\arrayrulewidth #1}\arrayrulecolor[HTML]{#2}\cline{#3}}

\begin{longtable}[c]{|p{1.20in}|p{1.20in}|p{1.20in}|p{1.20in}|p{1.20in}}

\caption{\textcolor[HTML]{0000AA}{\fontsize{14}{17}\selectfont{\textbf{Órgãos\ com\ maior\ percentual\ de\ indígenas}}}\textcolor[HTML]{0000AA}{\fontsize{14}{17}\selectfont{\textbf{\linebreak }}}\textcolor[HTML]{0000AA}{\fontsize{14}{17}\selectfont{\textbf{(6\ órgãos\ com\ maior\ representatividade)}}}}\\

\hhline{>{\arrayrulecolor[HTML]{666666}\global\arrayrulewidth=1.5pt}->{\arrayrulecolor[HTML]{666666}\global\arrayrulewidth=1.5pt}->{\arrayrulecolor[HTML]{666666}\global\arrayrulewidth=1.5pt}->{\arrayrulecolor[HTML]{666666}\global\arrayrulewidth=1.5pt}->{\arrayrulecolor[HTML]{666666}\global\arrayrulewidth=1.5pt}-}

\multicolumn{1}{>{\cellcolor[HTML]{0000AA}\raggedright}m{\dimexpr 1.2in+0\tabcolsep}}{\textcolor[HTML]{FFFFFF}{\fontsize{11}{11}\selectfont{Sigla\ do\ Órgão}}} & \multicolumn{1}{>{\cellcolor[HTML]{0000AA}\raggedright}m{\dimexpr 1.2in+0\tabcolsep}}{\textcolor[HTML]{FFFFFF}{\fontsize{11}{11}\selectfont{Órgão}}} & \multicolumn{1}{>{\cellcolor[HTML]{0000AA}\raggedleft}m{\dimexpr 1.2in+0\tabcolsep}}{\textcolor[HTML]{FFFFFF}{\fontsize{11}{11}\selectfont{Total\ de\ vínculos}}} & \multicolumn{1}{>{\cellcolor[HTML]{0000AA}\raggedleft}m{\dimexpr 1.2in+0\tabcolsep}}{\textcolor[HTML]{FFFFFF}{\fontsize{11}{11}\selectfont{Total\ de\ vínculos\ de\ indígenas}}} & \multicolumn{1}{>{\cellcolor[HTML]{0000AA}\raggedleft}m{\dimexpr 1.2in+0\tabcolsep}}{\textcolor[HTML]{FFFFFF}{\fontsize{11}{11}\selectfont{Participação\ (\%)\ de\ indígenas\ no\ órgão}}} \\

\noalign{\global\arrayrulewidth 0pt}\arrayrulecolor[HTML]{000000}

\hhline{>{\arrayrulecolor[HTML]{666666}\global\arrayrulewidth=1.5pt}->{\arrayrulecolor[HTML]{666666}\global\arrayrulewidth=1.5pt}->{\arrayrulecolor[HTML]{666666}\global\arrayrulewidth=1.5pt}->{\arrayrulecolor[HTML]{666666}\global\arrayrulewidth=1.5pt}->{\arrayrulecolor[HTML]{666666}\global\arrayrulewidth=1.5pt}-}\endfirsthead \caption[]{\textcolor[HTML]{0000AA}{\fontsize{14}{17}\selectfont{\textbf{Órgãos\ com\ maior\ percentual\ de\ indígenas}}}\textcolor[HTML]{0000AA}{\fontsize{14}{17}\selectfont{\textbf{\linebreak }}}\textcolor[HTML]{0000AA}{\fontsize{14}{17}\selectfont{\textbf{(6\ órgãos\ com\ maior\ representatividade)}}}}\\

\hhline{>{\arrayrulecolor[HTML]{666666}\global\arrayrulewidth=1.5pt}->{\arrayrulecolor[HTML]{666666}\global\arrayrulewidth=1.5pt}->{\arrayrulecolor[HTML]{666666}\global\arrayrulewidth=1.5pt}->{\arrayrulecolor[HTML]{666666}\global\arrayrulewidth=1.5pt}->{\arrayrulecolor[HTML]{666666}\global\arrayrulewidth=1.5pt}-}

\multicolumn{1}{>{\cellcolor[HTML]{0000AA}\raggedright}m{\dimexpr 1.2in+0\tabcolsep}}{\textcolor[HTML]{FFFFFF}{\fontsize{11}{11}\selectfont{Sigla\ do\ Órgão}}} & \multicolumn{1}{>{\cellcolor[HTML]{0000AA}\raggedright}m{\dimexpr 1.2in+0\tabcolsep}}{\textcolor[HTML]{FFFFFF}{\fontsize{11}{11}\selectfont{Órgão}}} & \multicolumn{1}{>{\cellcolor[HTML]{0000AA}\raggedleft}m{\dimexpr 1.2in+0\tabcolsep}}{\textcolor[HTML]{FFFFFF}{\fontsize{11}{11}\selectfont{Total\ de\ vínculos}}} & \multicolumn{1}{>{\cellcolor[HTML]{0000AA}\raggedleft}m{\dimexpr 1.2in+0\tabcolsep}}{\textcolor[HTML]{FFFFFF}{\fontsize{11}{11}\selectfont{Total\ de\ vínculos\ de\ indígenas}}} & \multicolumn{1}{>{\cellcolor[HTML]{0000AA}\raggedleft}m{\dimexpr 1.2in+0\tabcolsep}}{\textcolor[HTML]{FFFFFF}{\fontsize{11}{11}\selectfont{Participação\ (\%)\ de\ indígenas\ no\ órgão}}} \\

\noalign{\global\arrayrulewidth 0pt}\arrayrulecolor[HTML]{000000}

\hhline{>{\arrayrulecolor[HTML]{666666}\global\arrayrulewidth=1.5pt}->{\arrayrulecolor[HTML]{666666}\global\arrayrulewidth=1.5pt}->{\arrayrulecolor[HTML]{666666}\global\arrayrulewidth=1.5pt}->{\arrayrulecolor[HTML]{666666}\global\arrayrulewidth=1.5pt}->{\arrayrulecolor[HTML]{666666}\global\arrayrulewidth=1.5pt}-}\endhead



\multicolumn{1}{>{\raggedright}m{\dimexpr 1.2in+0\tabcolsep}}{\textcolor[HTML]{000000}{\fontsize{11}{11}\selectfont{MPI}}} & \multicolumn{1}{>{\raggedright}m{\dimexpr 1.2in+0\tabcolsep}}{\textcolor[HTML]{000000}{\fontsize{11}{11}\selectfont{Ministério\ dos\ Povos\ Indígenas}}} & \multicolumn{1}{>{\raggedleft}m{\dimexpr 1.2in+0\tabcolsep}}{\textcolor[HTML]{000000}{\fontsize{11}{11}\selectfont{109}}} & \multicolumn{1}{>{\raggedleft}m{\dimexpr 1.2in+0\tabcolsep}}{\textcolor[HTML]{000000}{\fontsize{11}{11}\selectfont{36}}} & \multicolumn{1}{>{\cellcolor[HTML]{00A100}\raggedleft}m{\dimexpr 1.2in+0\tabcolsep}}{\textcolor[HTML]{000000}{\fontsize{11}{11}\selectfont{33,0\%}}} \\

\noalign{\global\arrayrulewidth 0pt}\arrayrulecolor[HTML]{000000}





\multicolumn{1}{>{\raggedright}m{\dimexpr 1.2in+0\tabcolsep}}{\textcolor[HTML]{000000}{\fontsize{11}{11}\selectfont{FUNAI}}} & \multicolumn{1}{>{\raggedright}m{\dimexpr 1.2in+0\tabcolsep}}{\textcolor[HTML]{000000}{\fontsize{11}{11}\selectfont{Fundação\ Nacional\ dos\ Povos\ Indígenas}}} & \multicolumn{1}{>{\raggedleft}m{\dimexpr 1.2in+0\tabcolsep}}{\textcolor[HTML]{000000}{\fontsize{11}{11}\selectfont{2.047}}} & \multicolumn{1}{>{\raggedleft}m{\dimexpr 1.2in+0\tabcolsep}}{\textcolor[HTML]{000000}{\fontsize{11}{11}\selectfont{478}}} & \multicolumn{1}{>{\cellcolor[HTML]{6EBE5C}\raggedleft}m{\dimexpr 1.2in+0\tabcolsep}}{\textcolor[HTML]{000000}{\fontsize{11}{11}\selectfont{23,4\%}}} \\

\noalign{\global\arrayrulewidth 0pt}\arrayrulecolor[HTML]{000000}





\multicolumn{1}{>{\raggedright}m{\dimexpr 1.2in+0\tabcolsep}}{\textcolor[HTML]{000000}{\fontsize{11}{11}\selectfont{IBAMA}}} & \multicolumn{1}{>{\raggedright}m{\dimexpr 1.2in+0\tabcolsep}}{\textcolor[HTML]{000000}{\fontsize{11}{11}\selectfont{Instituto\ Brasileiro\ do\ Meio\ Ambiente\ e\ dos\ Recursos\ Naturais\ Renováveis}}} & \multicolumn{1}{>{\raggedleft}m{\dimexpr 1.2in+0\tabcolsep}}{\textcolor[HTML]{000000}{\fontsize{11}{11}\selectfont{3.722}}} & \multicolumn{1}{>{\raggedleft}m{\dimexpr 1.2in+0\tabcolsep}}{\textcolor[HTML]{000000}{\fontsize{11}{11}\selectfont{237}}} & \multicolumn{1}{>{\cellcolor[HTML]{D9EED2}\raggedleft}m{\dimexpr 1.2in+0\tabcolsep}}{\textcolor[HTML]{000000}{\fontsize{11}{11}\selectfont{6,4\%}}} \\

\noalign{\global\arrayrulewidth 0pt}\arrayrulecolor[HTML]{000000}





\multicolumn{1}{>{\raggedright}m{\dimexpr 1.2in+0\tabcolsep}}{\textcolor[HTML]{000000}{\fontsize{11}{11}\selectfont{IFAM}}} & \multicolumn{1}{>{\raggedright}m{\dimexpr 1.2in+0\tabcolsep}}{\textcolor[HTML]{000000}{\fontsize{11}{11}\selectfont{Instituto\ Federal\ do\ Amazonas}}} & \multicolumn{1}{>{\raggedleft}m{\dimexpr 1.2in+0\tabcolsep}}{\textcolor[HTML]{000000}{\fontsize{11}{11}\selectfont{1.983}}} & \multicolumn{1}{>{\raggedleft}m{\dimexpr 1.2in+0\tabcolsep}}{\textcolor[HTML]{000000}{\fontsize{11}{11}\selectfont{46}}} & \multicolumn{1}{>{\cellcolor[HTML]{F1F9EF}\raggedleft}m{\dimexpr 1.2in+0\tabcolsep}}{\textcolor[HTML]{000000}{\fontsize{11}{11}\selectfont{2,3\%}}} \\

\noalign{\global\arrayrulewidth 0pt}\arrayrulecolor[HTML]{000000}





\multicolumn{1}{>{\raggedright}m{\dimexpr 1.2in+0\tabcolsep}}{\textcolor[HTML]{000000}{\fontsize{11}{11}\selectfont{EX-TER/RR}}} & \multicolumn{1}{>{\raggedright}m{\dimexpr 1.2in+0\tabcolsep}}{\textcolor[HTML]{000000}{\fontsize{11}{11}\selectfont{Governo\ do\ Ex-Território\ de\ Roraima}}} & \multicolumn{1}{>{\raggedleft}m{\dimexpr 1.2in+0\tabcolsep}}{\textcolor[HTML]{000000}{\fontsize{11}{11}\selectfont{6.474}}} & \multicolumn{1}{>{\raggedleft}m{\dimexpr 1.2in+0\tabcolsep}}{\textcolor[HTML]{000000}{\fontsize{11}{11}\selectfont{128}}} & \multicolumn{1}{>{\cellcolor[HTML]{F3FAF1}\raggedleft}m{\dimexpr 1.2in+0\tabcolsep}}{\textcolor[HTML]{000000}{\fontsize{11}{11}\selectfont{2,0\%}}} \\

\noalign{\global\arrayrulewidth 0pt}\arrayrulecolor[HTML]{000000}





\multicolumn{1}{>{\raggedright}m{\dimexpr 1.2in+0\tabcolsep}}{\textcolor[HTML]{000000}{\fontsize{11}{11}\selectfont{UFRR}}} & \multicolumn{1}{>{\raggedright}m{\dimexpr 1.2in+0\tabcolsep}}{\textcolor[HTML]{000000}{\fontsize{11}{11}\selectfont{Universidade\ Federal\ de\ Roraima}}} & \multicolumn{1}{>{\raggedleft}m{\dimexpr 1.2in+0\tabcolsep}}{\textcolor[HTML]{000000}{\fontsize{11}{11}\selectfont{1.144}}} & \multicolumn{1}{>{\raggedleft}m{\dimexpr 1.2in+0\tabcolsep}}{\textcolor[HTML]{000000}{\fontsize{11}{11}\selectfont{21}}} & \multicolumn{1}{>{\cellcolor[HTML]{F4FAF2}\raggedleft}m{\dimexpr 1.2in+0\tabcolsep}}{\textcolor[HTML]{000000}{\fontsize{11}{11}\selectfont{1,8\%}}} \\

\noalign{\global\arrayrulewidth 0pt}\arrayrulecolor[HTML]{000000}

\hhline{>{\arrayrulecolor[HTML]{666666}\global\arrayrulewidth=1.5pt}->{\arrayrulecolor[HTML]{666666}\global\arrayrulewidth=1.5pt}->{\arrayrulecolor[HTML]{666666}\global\arrayrulewidth=1.5pt}->{\arrayrulecolor[HTML]{666666}\global\arrayrulewidth=1.5pt}->{\arrayrulecolor[HTML]{666666}\global\arrayrulewidth=1.5pt}-}



\end{longtable}



\arrayrulecolor[HTML]{000000}

\global\setlength{\arrayrulewidth}{\Oldarrayrulewidth}

\global\setlength{\tabcolsep}{\Oldtabcolsep}

\renewcommand*{\arraystretch}{1}

\section{Perfil dos vínculos ativos de
indígenas}\label{perfil-dos-vuxednculos-ativos-de-induxedgenas}

\subsection{Sexo e faixa etária}\label{sexo-e-faixa-etuxe1ria}

O conjunto de vínculos ativos de indígenas é composto por 33,17\% de
mulheres. Considerando todos os servidores da APF, este percentual é de
45,38\%

Quanto à idade, 34,51\% dos indígenas na APF têm até 40 anos e 8,59\%
estão acima dos 65. Estes percenutais em toda a APF são,
respectivamente, 27,43\% e 8,38\%

\pandocbounded{\includegraphics[keepaspectratio]{index_files/figure-latex/piramide_etaria-1.png}}
\pandocbounded{\includegraphics[keepaspectratio]{index_files/figure-latex/piramide_etaria-2.png}}

\subsection{Escolaridade}\label{escolaridade}

A figura a seguir apresenta a distribuição de vínculos que se
autodeclaram indígenas segundo a escolaridade. Cerca de 1.361 indígenas,
60,6\% do total, possuem nível superior ou maior. Considerando todos os
servidores da APF, este percentual é de 84,0\%.

\pandocbounded{\includegraphics[keepaspectratio]{index_files/figure-latex/treemap_indigena-1.png}}

\subsection{Órgão de Vínculo}\label{uxf3rguxe3o-de-vuxednculo}

Em Fevereiro/2026, a maior parte dos vínculos de indígenas, 39,8\%,
estavam em órgãos de Autarquia Federal. Essa categoria apresentou seu
menor percentual em 2016, com 33,0\% dos indígenas, e seu valor máximo
em 2024, com 48,2\%.

\pandocbounded{\includegraphics[keepaspectratio]{index_files/figure-latex/grafico_natjur-1.png}}

Ao analisar a distribuição de indígenas por órgão, a tabela a seguir
aponta Fundação Nacional Dos Povos Indígenas (FUNAI), Ministério Da
Saúde (MS), Instituto Brasileiro Do Meio Ambiente E Dos Recursos
Naturais Renováveis (IBAMA) e Governo Do Ex-Território De Roraima
(EX-TER/RR) como aqueles em que os vínculos de indígenas estão mais
concentrados, responsáveis por 49,5\% dos vínculos ativos de indígenas.

\global\setlength{\Oldarrayrulewidth}{\arrayrulewidth}

\global\setlength{\Oldtabcolsep}{\tabcolsep}

\setlength{\tabcolsep}{2pt}

\renewcommand*{\arraystretch}{1.5}



\providecommand{\ascline}[3]{\noalign{\global\arrayrulewidth #1}\arrayrulecolor[HTML]{#2}\cline{#3}}

\begin{longtable}[c]{|p{1.20in}|p{1.20in}|p{1.20in}|p{1.20in}|p{1.20in}}

\caption{\textcolor[HTML]{0000AA}{\fontsize{14}{17}\selectfont{\textbf{Percentual\ dos\ indígenas\ vinculados\ ao\ órgão}}}\textcolor[HTML]{0000AA}{\fontsize{14}{17}\selectfont{\textbf{\linebreak }}}\textcolor[HTML]{0000AA}{\fontsize{14}{17}\selectfont{\textbf{(6\ órgãos\ com\ maior\ percentual)}}}}\\

\hhline{>{\arrayrulecolor[HTML]{666666}\global\arrayrulewidth=1.5pt}->{\arrayrulecolor[HTML]{666666}\global\arrayrulewidth=1.5pt}->{\arrayrulecolor[HTML]{666666}\global\arrayrulewidth=1.5pt}->{\arrayrulecolor[HTML]{666666}\global\arrayrulewidth=1.5pt}->{\arrayrulecolor[HTML]{666666}\global\arrayrulewidth=1.5pt}-}

\multicolumn{1}{>{\cellcolor[HTML]{0000AA}\raggedright}m{\dimexpr 1.2in+0\tabcolsep}}{\textcolor[HTML]{FFFFFF}{\fontsize{11}{11}\selectfont{Sigla}}} & \multicolumn{1}{>{\cellcolor[HTML]{0000AA}\raggedright}m{\dimexpr 1.2in+0\tabcolsep}}{\textcolor[HTML]{FFFFFF}{\fontsize{11}{11}\selectfont{Órgão}}} & \multicolumn{1}{>{\cellcolor[HTML]{0000AA}\raggedleft}m{\dimexpr 1.2in+0\tabcolsep}}{\textcolor[HTML]{FFFFFF}{\fontsize{11}{11}\selectfont{Total\ de\ vínculos}}} & \multicolumn{1}{>{\cellcolor[HTML]{0000AA}\raggedleft}m{\dimexpr 1.2in+0\tabcolsep}}{\textcolor[HTML]{FFFFFF}{\fontsize{11}{11}\selectfont{Total\ de\ vínculos\ de\ indígenas}}} & \multicolumn{1}{>{\cellcolor[HTML]{0000AA}\raggedleft}m{\dimexpr 1.2in+0\tabcolsep}}{\textcolor[HTML]{FFFFFF}{\fontsize{11}{11}\selectfont{Percentual\ dos\ indígenas\ vinculados\ ao\ órgão}}} \\

\noalign{\global\arrayrulewidth 0pt}\arrayrulecolor[HTML]{000000}

\hhline{>{\arrayrulecolor[HTML]{666666}\global\arrayrulewidth=1.5pt}->{\arrayrulecolor[HTML]{666666}\global\arrayrulewidth=1.5pt}->{\arrayrulecolor[HTML]{666666}\global\arrayrulewidth=1.5pt}->{\arrayrulecolor[HTML]{666666}\global\arrayrulewidth=1.5pt}->{\arrayrulecolor[HTML]{666666}\global\arrayrulewidth=1.5pt}-}\endfirsthead \caption[]{\textcolor[HTML]{0000AA}{\fontsize{14}{17}\selectfont{\textbf{Percentual\ dos\ indígenas\ vinculados\ ao\ órgão}}}\textcolor[HTML]{0000AA}{\fontsize{14}{17}\selectfont{\textbf{\linebreak }}}\textcolor[HTML]{0000AA}{\fontsize{14}{17}\selectfont{\textbf{(6\ órgãos\ com\ maior\ percentual)}}}}\\

\hhline{>{\arrayrulecolor[HTML]{666666}\global\arrayrulewidth=1.5pt}->{\arrayrulecolor[HTML]{666666}\global\arrayrulewidth=1.5pt}->{\arrayrulecolor[HTML]{666666}\global\arrayrulewidth=1.5pt}->{\arrayrulecolor[HTML]{666666}\global\arrayrulewidth=1.5pt}->{\arrayrulecolor[HTML]{666666}\global\arrayrulewidth=1.5pt}-}

\multicolumn{1}{>{\cellcolor[HTML]{0000AA}\raggedright}m{\dimexpr 1.2in+0\tabcolsep}}{\textcolor[HTML]{FFFFFF}{\fontsize{11}{11}\selectfont{Sigla}}} & \multicolumn{1}{>{\cellcolor[HTML]{0000AA}\raggedright}m{\dimexpr 1.2in+0\tabcolsep}}{\textcolor[HTML]{FFFFFF}{\fontsize{11}{11}\selectfont{Órgão}}} & \multicolumn{1}{>{\cellcolor[HTML]{0000AA}\raggedleft}m{\dimexpr 1.2in+0\tabcolsep}}{\textcolor[HTML]{FFFFFF}{\fontsize{11}{11}\selectfont{Total\ de\ vínculos}}} & \multicolumn{1}{>{\cellcolor[HTML]{0000AA}\raggedleft}m{\dimexpr 1.2in+0\tabcolsep}}{\textcolor[HTML]{FFFFFF}{\fontsize{11}{11}\selectfont{Total\ de\ vínculos\ de\ indígenas}}} & \multicolumn{1}{>{\cellcolor[HTML]{0000AA}\raggedleft}m{\dimexpr 1.2in+0\tabcolsep}}{\textcolor[HTML]{FFFFFF}{\fontsize{11}{11}\selectfont{Percentual\ dos\ indígenas\ vinculados\ ao\ órgão}}} \\

\noalign{\global\arrayrulewidth 0pt}\arrayrulecolor[HTML]{000000}

\hhline{>{\arrayrulecolor[HTML]{666666}\global\arrayrulewidth=1.5pt}->{\arrayrulecolor[HTML]{666666}\global\arrayrulewidth=1.5pt}->{\arrayrulecolor[HTML]{666666}\global\arrayrulewidth=1.5pt}->{\arrayrulecolor[HTML]{666666}\global\arrayrulewidth=1.5pt}->{\arrayrulecolor[HTML]{666666}\global\arrayrulewidth=1.5pt}-}\endhead



\multicolumn{1}{>{\raggedright}m{\dimexpr 1.2in+0\tabcolsep}}{\textcolor[HTML]{000000}{\fontsize{11}{11}\selectfont{FUNAI}}} & \multicolumn{1}{>{\raggedright}m{\dimexpr 1.2in+0\tabcolsep}}{\textcolor[HTML]{000000}{\fontsize{11}{11}\selectfont{Fundação\ Nacional\ dos\ Povos\ Indígenas}}} & \multicolumn{1}{>{\raggedleft}m{\dimexpr 1.2in+0\tabcolsep}}{\textcolor[HTML]{000000}{\fontsize{11}{11}\selectfont{2.047}}} & \multicolumn{1}{>{\raggedleft}m{\dimexpr 1.2in+0\tabcolsep}}{\textcolor[HTML]{000000}{\fontsize{11}{11}\selectfont{478}}} & \multicolumn{1}{>{\cellcolor[HTML]{00A100}\raggedleft}m{\dimexpr 1.2in+0\tabcolsep}}{\textcolor[HTML]{000000}{\fontsize{11}{11}\selectfont{21,3\%}}} \\

\noalign{\global\arrayrulewidth 0pt}\arrayrulecolor[HTML]{000000}





\multicolumn{1}{>{\raggedright}m{\dimexpr 1.2in+0\tabcolsep}}{\textcolor[HTML]{000000}{\fontsize{11}{11}\selectfont{MS}}} & \multicolumn{1}{>{\raggedright}m{\dimexpr 1.2in+0\tabcolsep}}{\textcolor[HTML]{000000}{\fontsize{11}{11}\selectfont{Ministério\ da\ Saúde}}} & \multicolumn{1}{>{\raggedleft}m{\dimexpr 1.2in+0\tabcolsep}}{\textcolor[HTML]{000000}{\fontsize{11}{11}\selectfont{61.229}}} & \multicolumn{1}{>{\raggedleft}m{\dimexpr 1.2in+0\tabcolsep}}{\textcolor[HTML]{000000}{\fontsize{11}{11}\selectfont{268}}} & \multicolumn{1}{>{\cellcolor[HTML]{8FCC7D}\raggedleft}m{\dimexpr 1.2in+0\tabcolsep}}{\textcolor[HTML]{000000}{\fontsize{11}{11}\selectfont{11,9\%}}} \\

\noalign{\global\arrayrulewidth 0pt}\arrayrulecolor[HTML]{000000}





\multicolumn{1}{>{\raggedright}m{\dimexpr 1.2in+0\tabcolsep}}{\textcolor[HTML]{000000}{\fontsize{11}{11}\selectfont{IBAMA}}} & \multicolumn{1}{>{\raggedright}m{\dimexpr 1.2in+0\tabcolsep}}{\textcolor[HTML]{000000}{\fontsize{11}{11}\selectfont{Instituto\ Brasileiro\ do\ Meio\ Ambiente\ e\ dos\ Recursos\ Naturais\ Renováveis}}} & \multicolumn{1}{>{\raggedleft}m{\dimexpr 1.2in+0\tabcolsep}}{\textcolor[HTML]{000000}{\fontsize{11}{11}\selectfont{3.722}}} & \multicolumn{1}{>{\raggedleft}m{\dimexpr 1.2in+0\tabcolsep}}{\textcolor[HTML]{000000}{\fontsize{11}{11}\selectfont{237}}} & \multicolumn{1}{>{\cellcolor[HTML]{9CD28C}\raggedleft}m{\dimexpr 1.2in+0\tabcolsep}}{\textcolor[HTML]{000000}{\fontsize{11}{11}\selectfont{10,6\%}}} \\

\noalign{\global\arrayrulewidth 0pt}\arrayrulecolor[HTML]{000000}





\multicolumn{1}{>{\raggedright}m{\dimexpr 1.2in+0\tabcolsep}}{\textcolor[HTML]{000000}{\fontsize{11}{11}\selectfont{EX-TER/RR}}} & \multicolumn{1}{>{\raggedright}m{\dimexpr 1.2in+0\tabcolsep}}{\textcolor[HTML]{000000}{\fontsize{11}{11}\selectfont{Governo\ do\ Ex-Território\ de\ Roraima}}} & \multicolumn{1}{>{\raggedleft}m{\dimexpr 1.2in+0\tabcolsep}}{\textcolor[HTML]{000000}{\fontsize{11}{11}\selectfont{6.474}}} & \multicolumn{1}{>{\raggedleft}m{\dimexpr 1.2in+0\tabcolsep}}{\textcolor[HTML]{000000}{\fontsize{11}{11}\selectfont{128}}} & \multicolumn{1}{>{\cellcolor[HTML]{CBE7C0}\raggedleft}m{\dimexpr 1.2in+0\tabcolsep}}{\textcolor[HTML]{000000}{\fontsize{11}{11}\selectfont{5,7\%}}} \\

\noalign{\global\arrayrulewidth 0pt}\arrayrulecolor[HTML]{000000}





\multicolumn{1}{>{\raggedright}m{\dimexpr 1.2in+0\tabcolsep}}{\textcolor[HTML]{000000}{\fontsize{11}{11}\selectfont{ICMBIO}}} & \multicolumn{1}{>{\raggedright}m{\dimexpr 1.2in+0\tabcolsep}}{\textcolor[HTML]{000000}{\fontsize{11}{11}\selectfont{Instituto\ Chico\ Mendes\ de\ Conservação\ da\ Biodiversidade}}} & \multicolumn{1}{>{\raggedleft}m{\dimexpr 1.2in+0\tabcolsep}}{\textcolor[HTML]{000000}{\fontsize{11}{11}\selectfont{5.998}}} & \multicolumn{1}{>{\raggedleft}m{\dimexpr 1.2in+0\tabcolsep}}{\textcolor[HTML]{000000}{\fontsize{11}{11}\selectfont{98}}} & \multicolumn{1}{>{\cellcolor[HTML]{D7EDCF}\raggedleft}m{\dimexpr 1.2in+0\tabcolsep}}{\textcolor[HTML]{000000}{\fontsize{11}{11}\selectfont{4,4\%}}} \\

\noalign{\global\arrayrulewidth 0pt}\arrayrulecolor[HTML]{000000}





\multicolumn{1}{>{\raggedright}m{\dimexpr 1.2in+0\tabcolsep}}{\textcolor[HTML]{000000}{\fontsize{11}{11}\selectfont{IFAM}}} & \multicolumn{1}{>{\raggedright}m{\dimexpr 1.2in+0\tabcolsep}}{\textcolor[HTML]{000000}{\fontsize{11}{11}\selectfont{Instituto\ Federal\ do\ Amazonas}}} & \multicolumn{1}{>{\raggedleft}m{\dimexpr 1.2in+0\tabcolsep}}{\textcolor[HTML]{000000}{\fontsize{11}{11}\selectfont{1.983}}} & \multicolumn{1}{>{\raggedleft}m{\dimexpr 1.2in+0\tabcolsep}}{\textcolor[HTML]{000000}{\fontsize{11}{11}\selectfont{46}}} & \multicolumn{1}{>{\cellcolor[HTML]{ECF6E8}\raggedleft}m{\dimexpr 1.2in+0\tabcolsep}}{\textcolor[HTML]{000000}{\fontsize{11}{11}\selectfont{2,0\%}}} \\

\noalign{\global\arrayrulewidth 0pt}\arrayrulecolor[HTML]{000000}

\hhline{>{\arrayrulecolor[HTML]{666666}\global\arrayrulewidth=1.5pt}->{\arrayrulecolor[HTML]{666666}\global\arrayrulewidth=1.5pt}->{\arrayrulecolor[HTML]{666666}\global\arrayrulewidth=1.5pt}->{\arrayrulecolor[HTML]{666666}\global\arrayrulewidth=1.5pt}->{\arrayrulecolor[HTML]{666666}\global\arrayrulewidth=1.5pt}-}



\end{longtable}



\arrayrulecolor[HTML]{000000}

\global\setlength{\arrayrulewidth}{\Oldarrayrulewidth}

\global\setlength{\tabcolsep}{\Oldtabcolsep}

\renewcommand*{\arraystretch}{1}

Com relação à UF de exercício do servidor, a maior concentração de
vínculos indígena está no DF, RR e AM, que somam 966, 43,0\% do total
deste grupo.

\pandocbounded{\includegraphics[keepaspectratio]{index_files/figure-latex/show_mapa-1.png}}

\section{Saídas e ingressos de servidores
efetivos}\label{sauxeddas-e-ingressos-de-servidores-efetivos}

\subsection{Servidores efetivos}\label{servidores-efetivos}

Em Fevereiro/2026, 73\% dos vínculos ativos de indígenas eram de
servidores com cargo efetivo. Considerando todos os vínculos ativos da
APF, este percentual era de 94\%. Entre os indígenas, o percentuaL de
servidores com cargo efetivo era de 75.3\%.

\pandocbounded{\includegraphics[keepaspectratio]{index_files/figure-latex/plot_efetivos-1.png}}

Ao longo do período de 2016 a 2026, o total de vínculos efetivos de
indígenas variou de forma considerável. A série começa em 2016 com 1.844
indígenas efetivos, chegando ao seu menor valor, 1.453 em 2022. Tal
diminuição é explicada por aposentadorias e óbitos no período, sem a
contrapartida de ingressos indígenas suficientes à manutenção do
quantitativo. Todavia, a partir de 2022, o quantitativo de indígenas vem
apresentando aumento.

Em relação aos vínculos sem cargo efetivo, também registrou queda a
partir de 2016, com um salto no ano 2021 de 416 para 800.

\pandocbounded{\includegraphics[keepaspectratio]{index_files/figure-latex/indigenas_compara_efet-1.png}}

Tais evoluções impactaram na evolução da proporção de vínculos efetivos
de indígenas. Em 2016, 80,5\% dos vínculos ativos de indígenas eram
efetivos. Este percentual sofre a maior queda de 2020 para 2021
(coincidindo com o salto no número de não efetivos), chegando a 65,5\%,
alcançando seu menor valor de 62,8\% em 2023, finalizando em 2026 com
75,3\%.

Dentre os não-indígenas, a maior variação do percentual de vínculos
efetivos ocorreu de 2016 para 2017, saltando de 80.5\% para 81.4\%.

\pandocbounded{\includegraphics[keepaspectratio]{index_files/figure-latex/plot_prop_efetivos-1.png}}

\subsection{Saídas e ingressos (com cotas x sem
cotas)}\label{sauxeddas-e-ingressos-com-cotas-x-sem-cotas}

O gráfico a seguir aponta as variações anuais de saídas e ingressos de
servidores efetivos indígenas. A variação no número de efetivos foi
negativa durante todo o período em que o número de aposentadorias e
óbito foi superior ao número de ingressos, de 2017 a 2022. Tal situação
explica a queda a queda no número de servidores efetivos no período.

\pandocbounded{\includegraphics[keepaspectratio]{index_files/figure-latex/variacao_efetivos-1.png}}

A reserva de vagas específica para indígenas em concursos públicos
federais foi estabelecida de forma abrangente com a Lei nº 15.142, de 3
de junho de 2025. Esta nova legislação substituiu a lei de cotas
anterior, de 2014, ampliando para 30\% a reserva de vagas para
candidatos pretos, pardos, indígenas e quilombolas. O percentual
estabelecido especificamente para indígenas foi de 3\%.

Desde então, foram registrados 224 ingressos por cota para indígenas, o
que corresponde a 77,8\% dos ingressos de indígenas (288) no governo
federal.

\pandocbounded{\includegraphics[keepaspectratio]{index_files/figure-latex/evolucao_efetivos-1.png}}

\section{Cargos Comissionados e Funções de
Confiança}\label{cargos-comissionados-e-funuxe7uxf5es-de-confianuxe7a}

Em Fevereiro/2026, a incidência de Cargos Comissionados e Funções de
confiança entre vínculos indígenas é de 19,0\%, superior à observada nas
demais categorias de raça/cor, na ordem de 16,9\%. As categorias onde se
observa maior parcela de vínculos com função são INDIGENA (19\%), BRANCA
(17.6\%) e PARDA (16.2\%).

Quanto ao gênero, o percentual de exercício de funções entre homens
indígenas é 2,7 pontos percentuais acima do observado entre mulheres
indígenas. Nas demais categorias de cor/raça/etnia, a diferença é de 1,8
a mais entre homens em comparação às mulheres.

\pandocbounded{\includegraphics[keepaspectratio]{index_files/figure-latex/plot_funcao-1.png}}

O percentual de vínculos de indígenas que assumem alguma função observou
o seu valor mais baixo em dezembro de 2021, com 10.4\%, e seu maior
valor em fevereiro de 2026, com 19\%. O útimo período em que o
percentual de funções entre indígenas era menor que o observado para as
demais categorias foi em dezembro de 2022, quando 10.6\% dos indígenas
exerciam funções, enquanto este percentual nas demais categorias era de
15.8\%.

\pandocbounded{\includegraphics[keepaspectratio]{index_files/figure-latex/plot_serie_funcao-1.png}}

\subsection{Índice de equidade para função
FCE}\label{uxedndice-de-equidade-para-funuxe7uxe3o-fce}

Uma forma de analisar o acesso a níveis de liderança pode ser
determinada por um \emph{índice de equidade de acessos}. O indicador é
construído partindo-se da hipótese de que, na ausência de desigualdade
racial e de gênero no acesso à Função Comissionada Executiva (FCE), a
proporção de indígenas em cargos FCE deveria ser equivalente à sua
proporção entre os servidores efetivos ativos. Isso porque, para assumir
uma função FCE é necessário que o servidor tenha um vínculo ativo com
cargo efetivo. Assim, este índice é definido como sendo a razão entre o
percentual de indígenas em função FCE e o percentual de pessoas
indígenas no quadro efetivo ativo. Mais precisamente,

\[\text{Índice de equidade} = \frac{\text{% indígenas em cargos FCE}}{\text{% indígenas no quadro efetivo ativo}}
\]

O valor é interpretado de acordo com as seguintes faixas de valores:

\begin{itemize}
\item
  \(= 1\) → Equidade na ocupação dos cargos
\item
  \(< 1\) → Sub-representação de indígenas na liderança
\item
  \(> 1\) → Sobre-representação
\end{itemize}

O indicador pode ser desagregado por gênero e por nível de função. É
importante ressaltar que o indicador considera apenas os vínculos
efetivos e níveis de funções FCE.

Ao analisar a série histórica do indicador de equidade por nível FCE e
gênero, nota-se que os níveis com maior equidade variaram entre as
faixas 5 e 6, 13 e 14 e 15 a 18. Por outro lado, entre os homens, apenas
os níveis 5 e 6 apresentaram maior equidade.

\pandocbounded{\includegraphics[keepaspectratio]{index_files/figure-latex/razao_equivalencia_fce-1.png}}

\end{document}
